\subsection{Справочник <<Оборудование>>}
\label{spr:eqpt}


\subsubsection{Описание предметной области}

Справочник <<Оборудование>> содержит информацию о технологическом оборудовании на предприятии.

Справочник уже существует в СИСТЕМЕ, требуется внести изменения в его работу.

\subsubsection{Атрибуты}

В справочнике должны быть добавлены следующие поля.

\pc
\begin{longtable}{|p{50mm}|p{30mm}|p{78mm}|}
\hline
{\bf Наименование} & {\bf Тип данных} &  {\bf Комментарий} \endhead
\hline 
   Обмен с гофроагрегатом & Булево & Признак использования Обмена с гофроагрегатом \\
   \hline
   ТипРилевокПоУмолчанию & Перечисление. ТипыРилевок & Тип Рилевок По Умолчанию \\
   \hline
   ИнтервалОпросаБД & Число (10) & Интервал опроса БД, сек \\
   \hline
   СтрокаСоединенияСБД & Строка(0) & Строка соединения с БД системы ACS \\
   \hline
\caption{Описание свойств объекта <<Оборудование>>.}
\label{tab:tdb_eqpt}
\end{longtable} 


% \item Установлен флаг «Обмен с гофроагрегатом» – влияет на доступность команд обмена в документе выработки;
% \item ID Пустого контрагента – идентификатор контрагента, который будет подставляться для пустого заказа в раскрое (см. «Общие правила формирования заданий», п.3.2);
% \item Интервал опроса БД, сек. – частота обновления данных из базы JSDCS;



\subsubsection{Функциональные требования}

\point{Изменить форму ''Оборудование''} 

Добавить вкладку Обмен ACS. 
На вкладке расположить новые поля таблицы \ref{tab:tdb_eqpt}.

Строка подключения с БД – параметры подключения к базе asc\_dev, где:
 \begin{itemize}
     \item DRIVER –-- тип СУБД (MS SQL, MySQL и т.д.);
     \item SERVER --– адрес базы данных ACS в сети предприятия;
     \item ID –-- имя пользователя (По умолчанию  для asc\_dev ''sa'');
     \item PWD –-- пароль (По умолчанию  для asc\_dev ''sa'');
     \item DATABASE –-- имя базы данных (По умолчанию для asc\_dev ''asc\_dev'');
\end{itemize}

Добавить на вкладке Обмен ACS новую вкладку ''Проверка работы с JSM''



\point{Проверить соединение с JSM}

Добавить команду ''Проверить соединение с JSM''.
По команде пользователя необходимо выполнить тестовое подключение к базе данных JSM.

В случае неуспешного подключения необходимо выдавать сообщение об ошибке подключения.



\point{Получить данные с JSM}

По команде ''Получить данные с JSM'' система должна получить из памяти JSM текущие данные по выработки 
гофроагрегата. 

Результаты запроса должны быть выведены в отдельное поле на форме ''Результат запроса''.